\chapter{Background and Fundamentals}

This chapter reviews the theoretical and statistical principles foundational to this thesis.

\section{Single-Cell RNA Sequencing (scRNA-seq)}

\subsection{Overview}
Bulk RNA-seq averages gene expression across thousands of cells. Single-cell RNA sequencing (scRNA-seq) profiles each cell individually[cite: 1]. This reveals cellular heterogeneity hidden in bulk measurements. The output of an scRNA-seq workflow is an expression count matrix $X \in \mathbb{N}^{G \times C}$, where rows index genomic features (genes or transcripts), columns represent distinct single cells, and entry $x_{ij}$ records the observed read count for feature $i$ in cell $j$.

\subsection{Pseudo-Alignment}
Quantifying transcript expression requires mapping short sequencing reads to a reference transcriptome. Traditional alignment algorithms map reads base-by-base against a reference genome, a process that is computationally intensive[cite: 1]. Modern transcriptomics relies on \textbf{pseudo-alignment}: rather than finding exact base-level coordinates, pseudo-alignment algorithms (such as Salmon or alevin-fry) use index structures of length-$k$ substrings ($k$-mers) to identify the equivalence class of transcripts compatible with a given read[cite: 1]. By bypassing explicit coordinate mapping, pseudo-aligners achieve orders-of-magnitude speedups while maintaining high quantification accuracy.

\subsection{The Expectation-Maximization (EM) Algorithm}
A major challenge in pseudo-alignment is \textbf{multi-mapping reads}---reads that match multiple distinct transcript isoforms due to sequence similarity[cite: 1]. This problem is particularly acute for $3'$-biased or $5'$-biased protocols. When a read captures only the terminal region of a transcript, it may map equally well to all isoforms sharing that exon [cite: 1]. Full-length sequencing (e.g., SMART-seq2) often produces reads spanning unique internal exons that unambiguously identify the source transcript[cite: 1].

Pseudo-alignment algorithms like Salmon and alevin-fry resolve this ambiguity using the Expectation-Maximization (EM) algorithm. Rather than discarding multi-mapping reads, EM probabilistically assigns them to transcripts without losing information.

The EM algorithm is an iterative optimization technique used to find maximum likelihood estimates in statistical models with latent variables:
\begin{itemize}
    \item \textbf{Expectation (E-step):} For each multi-mapping read, calculate the probability it belongs to every candidate transcript.
    \item \textbf{Maximization (M-step):} Update global transcript abundances using these probabilities.
\end{itemize}

Each iteration refines abundance estimates. This assigns fractional read counts across isoforms while preserving relative signals[cite: 1]. For DTU detection, this matters—discarding ambiguous reads would lose the signal needed to distinguish similar isoforms.

\section{The Matrix Sparsity Problem}
Single-cell count matrices are characterized by extreme sparsity, with zero values typically accounting for 70--90\% of the matrix entries[cite: 1]. In count matrices, these zeros represent a mixture of two distinct phenomena:
\begin{itemize}
    \item \textbf{Biological Zeros:} True absences where a gene or transcript is not transcribed in a specific cell state.
    \item \textbf{Technical Zeros (Dropouts):} False absences where an mRNA molecule was present in the cell but failed to be captured, reverse-transcribed, or sequenced due to low chemical capture efficiency (typically 10--20\%).
\end{itemize}

Because biological and technical zeros are numerically indistinguishable in raw count matrices, sparsity introduces severe statistical hurdles.

\section{Limitations of Traditional Analytical Approaches}
Statistical methods designed for bulk RNA-seq (where each sample represents thousands of cells) assume relatively high read depth and low sparsity. When directly applied to zero-inflated single-cell count matrices, their mathematical assumptions are violated[cite: 1]:

\begin{enumerate}
    \item \textbf{Variance-Stabilizing Transformations:} Bulk analysis commonly uses $\log(x+c)$ (where $c$ is a pseudo-count) to normalize skewed count distributions. On zero-inflated matrices, adding an arbitrary pseudo-count and logging non-linearly distorts low-abundance measurements, compressing zeros into artificial scaling artifacts[cite: 1].
    
    \item \textbf{Linear Correlation Metrics:} Pearson and Spearman correlations assume smooth, continuous distributions. In sparse matrices, these metrics are distorted by the joint occurrence of technical zeros, measuring dropout coincidence rather than true regulatory co-expression[cite: 1].
    
    \item \textbf{Imputation Methods:} Heuristic smoothing replaces zeros with estimated values based on neighboring cells, frequently introducing artificial correlations and spurious downstream signals[cite: 1].
\end{enumerate}

Consequently, analyzing sparse single-cell data requires statistical models that treat zero observations explicitly rather than applying transformations designed for dense bulk data.

\section{Prior Work on Differential Transcript Usage (DTU)}
Differential Transcript Usage (DTU) analysis detects changes in the relative proportions of transcript isoforms between conditions or cell types, independent of overall gene expression[cite: 1]. Established bulk RNA-seq methods (such as rMATS, SUPPA2, or DRIMSeq) rely on continuous proportion modeling, junction-spanning read depths, or Dirichlet-multinomial distributions.

When applied to single-cell data, these methods largely fail because low per-cell sequencing depth results in degenerate proportion estimates, while terminal positional biases render internal splicing events invisible[cite: 1]. Furthermore, aggregating cells into pseudo-bulk profiles to recover read depth obscures the cell-to-cell heterogeneity that drives DTU.

\section{Statistical Foundations: Contingency Tables}
To handle extreme sparsity without imputation or transformation, the COTAN framework bypasses continuous count amplitudes and instead evaluates pairwise dependencies using asymptotic independence testing on binary presence/absence states[cite: 1].

For any two features A and B, their co-occurrence across $N$ cells is discretized into a $2\times2$ contingency table. Under the null hypothesis ($H_0$) of statistical independence, the expected joint frequency $E_{ij}$ is defined as:
$$E_{ij} = \frac{R_i \cdot C_j}{N}$$
where $R_i$ is the row marginal total, $C_j$ is the column marginal total, and $N$ is the total number of cells[cite: 1].

The deviation between observed counts $O_{ij}$ and expected counts $E_{ij}$ is evaluated via Pearson's chi-squared test statistic[cite: 1]:
$$\chi^2 = \sum_{i\in\{0,1\}} \sum_{j\in\{0,1\}} \frac{(O_{ij} - E_{ij})^2}{E_{ij}}$$

Because droplet scRNA-seq datasets typically contain thousands of cells ($N \gg 30$), the asymptotic chi-squared approximation remains robust[cite: 1]. This enables computationally efficient, large-scale hypothesis testing across millions of feature pairs without the bottleneck of exact inference methods.

\section{The COTAN Analytical Framework}
The COTAN (Co-expression Tables Analysis) framework is designed to model zero-inflated distributions directly without relying on data imputation or variance-stabilizing transformations[cite: 1].

The core premise of COTAN is that the binary occurrence pattern of features across thousands of cells carries structural regulatory signals. By discretizing expression into binary presence/absence states and evaluating pairwise associations using contingency-table statistics, COTAN is robust to sequencing depth variations and library size imbalances[cite: 1]. Within this framework, negative statistical associations (mutual exclusivity) between isoforms of the same gene serve as a primary indicator of alternative transcript usage.

COTAN generates three primary metrics[cite: 1]:
\begin{enumerate}
    \item \textbf{Co-expression Matrix (coex):} Measures normalized pairwise associations derived from contingency table statistics.
    
    \item \textbf{Differential Expression Analysis (DEA):} Quantifies feature detection probabilities across cellular subpopulations, enabling identification of cluster-specific markers.
    
    \item \textbf{Global Differentiation Index (GDI):} Identifies condition-specific markers by measuring how specifically a feature is expressed relative to the rest of the dataset; also serves as a homogeneity check for clustering quality.
\end{enumerate}

In summary, because current DTU algorithms fail on droplet scRNA-seq data due to matrix sparsity and positional bias, this thesis adapts the COTAN framework to transcript-level matrices to identify mutually exclusive isoform pairs without requiring artificial imputation or gene-level aggregation[cite: 1]. The combination of pseudo-alignment for transcript quantification, contingency-table statistics for sparse data modeling, and mutual exclusivity as a DTU indicator provides a theoretically grounded approach to recovering isoform switches from challenging single-cell datasets.
